% TOP_PROBLEM: {"problemId": "problem.zilber-s-quasiminimality-conjecture-for-the-complex-exponential-field", "canonicalTitle": "Zilber's Quasiminimality Conjecture for the Complex Exponential Field", "releaseRank": 475, "primaryDomain": "logic_foundations_set_theory", "primaryDomainLabel": "Logic, foundations & set theory", "displayStatus": "open", "releaseStatus": "open"}
% !TeX program = lualatex
% Definition and sources only; this file is not an LLM solution attempt.
\documentclass[11pt]{article}
\usepackage[margin=25mm]{geometry}
\usepackage{fontspec}
\setmainfont{Latin Modern Roman}
\usepackage{amsmath,amssymb,mathtools}
\usepackage{unicode-math}
\setmathfont{Latin Modern Math}
\usepackage{xurl,hyperref}
\hypersetup{colorlinks=true,urlcolor=blue}
\setcounter{secnumdepth}{0}
\setlength{\parindent}{0pt}
\setlength{\parskip}{5pt}
\setlength{\emergencystretch}{3em}
\begin{document}
\section*{\texorpdfstring{Zilber's Quasiminimality Conjecture for the Complex Exponential Field}{Zilber's Quasiminimality Conjecture for the Complex Exponential Field}}\label{475}
\subsection{Definitions and mathematical statement}
The complex exponential field is the first-order structure
\[
 {\mathbb{C}}_{\exp}=({\mathbb{C}};+ ,\cdot,0,1,\exp),\qquad \exp(z)=e^z.
\]
A subset $D\subseteq{\mathbb{C}}$ is definable with parameters if some first-order formula $\varphi(x,a_1,\ldots,a_m)$ in this language, with finitely many parameters from ${\mathbb{C}}$, satisfies $D=\{z:{\mathbb{C}}_{\exp}\models\varphi(z,a_1,\ldots,a_m)\}$. Quantifiers range over all complex numbers.

Quasiminimality asserts that every such one-variable definable set is countable or has countable complement. Countable includes finite. This is not the stronger claim that every definable set is finite or cofinite, nor a claim about arbitrary analytic subsets or arbitrary subsets of ${\mathbb{C}}$. Parameters are allowed, which is part of the requested scope.
\par\smallskip\textit{Formulation and concept references:} \href{https://arxiv.org/pdf/2306.14562v1}{[S1]}.

\subsection{Short English statement}
Is every one-variable set definable using complex arithmetic and exponentiation either countable or the complement of a countable set?
\subsection{Sources}
\begin{itemize}
\item[S1]A. J. Wilkie, "Analytic continuation and Zilber's quasiminimality conjecture," arXiv:2306.14562 (26 June 2023).
\url{https://arxiv.org/pdf/2306.14562v1}.
\textit{Formulation source.}
\end{itemize}
\end{document}
